\clearpage
\section*{Przedmowa}

Odkąd jeżdżę samochodem po Warszawie, zastanawiam się, jak faktycznie działa sygnalizacja świetlna. Na drogach wielokrotnie spotykałem się z sytuacjami, w których wydawało mi się, że sygnalizacja działa w sposób nieoptymalny. Jest to problem, który do dzisiaj nie daje mi spokoju.

Dlatego kiedy na studiach zainteresowałem się uczeniem maszynowym, zacząłem się zastanawiać, czy nie można by wykorzystać modeli uczenia ze wzmocnieniem właśnie do optymalizacji ruchu drogowego. Wydaje mi się, że jest to technologia z dużym potencjałem w tym zastosowaniu, ponieważ algorytm uczy się dostosowywać swoją politykę do konkretnego skrzyżowania. Nie są potrzebne ręcznie oznaczone zbiory treningowe, potrzebne są natomiast metoda obserwacji środowiska, możliwość podejmowania w nim akcji oraz funkcja nagrody.

Ponieważ kończę już studia na tym kierunku, chciałbym w tym miejscu podziękować moim kolegom i koleżankom za przejście tej drogi razem ze mną, wspólne projekty, naukę do kolokwiów i wzajemne motywowanie się do ukończenia studiów oraz kontynuowania nauki na studiach magisterskich. Bez Was na pewno byłoby dużo trudniej i nie byłaby to taka dobra zabawa.
